\begin{thebibliography}{200}  

    \bibitem{ref1}Although Japanese naming conventions place the family name first (e.g., Kaname Madoka), the official translation of Madoka Magica consistently uses the English convention of given name first, then family name. To avoid confusion, this book follows the order used in the show. For consistency, the English order will be maintained even when referring to characters in other works or real people.
    \bibitem{ref2}"The objects of imitation [in poetry] are men in action." Aristotle, Poetics (trans. S. H. Butcher) (The Internet Classics Archive, 2009). \url{http://classics.mit.edu/Aristotle/poetics.html}. Note that in Part 1 of Poetics Aristotle establishes that he includes drama and narrative as a part of poetry, and the way in which he uses "object of imitation" is what we would today call the subject of a work.
    \bibitem{ref3}Vivien Burr, "Where do you get your personality from?" An Introduction to Social Constructionism (London: Routledge, 1995).
    \bibitem{ref4}Ross Murfin and Supryia M. Ray, "Modernism," The Bedford Glossary of Critical and Literary Terms (Boston: Bedford Books, 1997).
    \bibitem{ref5}Murfin and Ray, "Postmodernism," The Bedford Glossary of Critical and Literary Terms. Note however that Murfin and Ray accentuate the negative and experimental aspects in their definition, and downplay the same pluralism which Hassan regards as a central feature.
    \bibitem{ref6}Ihab Hassan, "Pluralism in Postmodernism," Exploring Postmodernism: Selected Papers Presented at a Workshop on Postmodernism at the XIth International Comparative Literature Congress, Paris, 20–24 August 1985, Mattei Calinesku and Douew W. Fokkema ed. (Amsterdam: John Benjamins Publishing Company, 1987). Hassan provides a list of eleven overlapping criteria; the definition I provide draws heavily on his "Constructionism," "Decanonization," "Irony," "Carnivalization," and "Hybridization" criteria.
    \bibitem{ref7}Ibid.
    \bibitem{ref8}Pauline Marie Rosenau, "Affirmatives and Skeptics," The Fontana Postmodernism Reader, Walter Truett Anderson (ed.) (Fontana Press, 1996).
    \bibitem{ref9}Bill Shipley, "Vaclav Havel: Ambassador of Conscience 2003: From Prisoner to President—A Tribute," Art for Amnesty (Amnesty International, 2003). Accessed via Archive Today, \url{http://archive.today/B9l16}
    \bibitem{ref10}There is some ambiguity early in the film, which looks like it could be a very odd "alternate timeline" similar to the ones explored in the manga discussed in Part Two of this book, but marketing materials prior to the release of the movies referred to two recaps and a "continuation" of the Madoka Magica story. In the film itself, several characters are able to remember events from the final episode (which takes place after the point at which Homura reset all prior timelines, the fight with Walpurgisnacht), such as Mami and Homura talking about Wraiths, Kyubey remembering Homura telling him about Madoka's sacrifice to create a new system, and Madoka appearing at the end in "Madokanon" form.
    \bibitem{ref11}The chat in question is archived on my blog, JedABlue.com, at  \url{http://jedablue.com/2014/06/25/the-nutcracker-the-mouse-king-and-the-puella-magi/}
    \bibitem{ref12}Thomas Adajian, "The Definition of Art," Stanford Encyclopedia of Philosophy, Edward N. Zalta (ed.) \url{http://plato.stanford.edu/entries/art-definition/}
    \bibitem{ref13}Eric P. Nash, "Raising the Curtain," Manga Kamishibai (New York: Abrams Comicarts, 2009).
    \bibitem{ref14}Ibid.
    \bibitem{ref15}"Sakura and the Mysterious Magic Book," Cardcaptor Sakura: The Complete Collection (NIS America, 2014).
    \bibitem{ref16}"The Crybaby: Usagi's Beautiful Transformation," Pretty Soldier Sailor Moon (Vis Media, 2014).
    \bibitem{ref17}"The Black Claw Grips the Heart," Cutie Honey: The Complete Series (Discotek Media, 2013).
    \bibitem{ref18}"Take Care of Yourself," Neon Genesis Evangelion: Platinum Complete (ADV Films, 2005).
    \bibitem{ref19}Nagaru Tanigawa, The Melancholy of Haruhi Suzumiya (New York: Little, Brown and Company, 2009).
    \bibitem{ref20}"The Crybaby."
    \bibitem{ref21}Hassan, "Pluralism in Postmodernism."
    \bibitem{ref22}Brian Boyd, "Laughter and Literature: A Play Theory of Humor," Philosophy and Literature, Volume 28, Number 1 (April 2004).
    \bibitem{ref23}Noël Carroll, "Horror and Humor," The Journal of Aesthetics and Art Criticism, Volume 57, Number 2 (Spring 1999).
    \bibitem{ref24}Tomaki Saito, Beautiful Fighting Girl, J. Keith Vincent and Dawn Lawson (trans.) (Minneapolis, Minn.: University of Minneapolis Press, 2011)
    \bibitem{ref25}Sharalyn Orbaugh, "Busty Battlin' Babes: The Evolution of the Shojo in 1990s Visual Culture," Gender and Power in the Japanese Visual Field, Joshua S. Mostow, Norman Bryson, and Maribeth Graybill (ed.) (Honolulu: University of Hawai'i Press, 2003).
    \bibitem{ref26}Ibid. This concept of "Male Gaze" refers to a tendency in the framing of shots in film and television such that the camera emulates the prurient gaze of a heterosexual male viewer, for example by lingering on a woman's curves. See Laura Mulvey, "Visual Pleasure and Narrative Cinema," Media and Cultural Studies: Key Works, Meenakshi Gigi Durham and Douglas M. Kellner (ed.) (Malden, MA: Blackwell, 2006).
    \bibitem{ref27}Briefly, hegemonic masculinity refers to a cultural tendency to equate masculinity and power. Expressions of power such as violence, dominance, and assertiveness are seen as inherently or appropriately masculine. By contrast, performed femininity refers to a cultural tendency to equate femininity to a performance, a sort of social game by which women can create a constrained, but relatively safe, space for self-expression at the price of playing to the assumptions of hegemonic masculinity. See R. W. Connell, Gender and Power: Society, the Person and Sexual Politics (Stanford, CA: Stanford University Press, 1987).
    \bibitem{ref28}See for instance Anime News Network's dismissal of Magical Girl Lyrical Nanoha as "otaku-targeted entertainment taken to its logical extreme," noting in particular its "moe character types... and cheesecake service shots." Carl Kimlinger, "Magical Girl Lyrical Nanoha DVD Box Set: Review," Anime News Network (December 15, 2008). \url{http://www.animenewsnetwork.com/review/magical-girl-lyrical-nanoha/dvd-box-set}
    \bibitem{ref29}Saito, Beautiful Fighting Girl.
    \bibitem{ref30}"Sakura and the Mysterious Magic Book."
    \bibitem{ref31}Connell, Gender and Power.
    \bibitem{ref32}Namely, like Madoka, at the beginning of her second episode Sakura wakes in her bed just as she did in the beginning of the first episode, and initially assumes the events of the prior episode were a dream—until she realizes that one of the stuffed animals by her bed is not a stuffed animal at all. "Sakura's Wonderful Friend," \emph{Cardcaptor Sakura}.
    \bibitem{ref33}For a brief description of the Buddhist concept of impermanence designed to be readily accessible for Westerners more familiar with Christianity, see Ninian Smart, \emph{Buddhism and Christianity: Rivals and Allies} (Honolulu: University of Hawaii Press, 1993).
    \bibitem{ref34}Magica Quartet, Puella Magi Madoka Magica Production Note (SHAFT, 2011). The book is available only in Japanese; however, a rough translation of the relevant page is available at the unofficial Puella Magi Wiki: \url{http://wiki.puella-magi.net/Puella_Magi_Production_Note\#Witch_of_cheesecake_.28early_concept_that_later_became_Charlotte.29}.
    \bibitem{ref35}Notably, this is a stereotypically masculine-gendered activity, emphasizing the gender-bending aspect of Junko's characterization. See Romit Dasgupta, "Working with homosociality," \emph{The Salaryman in Japan: Crafting} Masculinities (Abingdon, Oxon: Routledge, 2013).
    \bibitem{ref36}"For putrefaction is the change and death of all things, and destruction of the first essence of all Natural things; whence there ariseth a regeneration, and new generation a thousand times better." Paracelsus, On the Nature of Things, excerpted in Stanton J. Linden (ed.), \emph{The Alchemy Reader: From Hermes Trismegistus to Isaac Newton} (Cambridge: Cambridge University Press, 2003).
    \bibitem{ref37}Gen Urobuchi, Fate/zero Volume 1 (Type-Moon, 2006). The book is available only in Japanese; an unauthorized translation of the relevant section is available at the pirate translation site Baka-Tsuki: \url{http://www.baka-tsuki.org/project/index.php?title=Fate/Zero:Volume_1_Postface_1}.
    \bibitem{ref38}Wayne C. Booth, The Rhetoric of Fiction (Chicago: University of Chicago Press, 1961).
    \bibitem{ref39}"The Body," Buffy the Vampire Slayer: The Complete Series (Fox Searchlight, 2010).
    \bibitem{ref40}Karen A. Smyers, "Symbolizing Inari: The Jewel," The Fox and the Jewel (Honolulu: University of Hawai'i Press, 1999).
    \bibitem{ref41}Princess Tutu: Complete Collection (Æsir Holdings, 2011).
    \bibitem{ref42}"The Spinners," Princess Tutu.
    \bibitem{ref43}Urobuchi, Fate/zero Volume 1.
    \bibitem{ref44}See for example Frazer's version of the tale. Sir James George Frazer, "The External Soul in Folk-Tales," The Golden Bough: A Study in Magic and Religion (Franklin Center, PA: Franklin Library, 1982)
    \bibitem{ref45}Jennifer Clarke Wilkes and Jon Pickens (ed.), "Lich," Dungeons \& Dragons: Monster Manual: Core Rulebook III v.3.5 (Renton, WA: Wizards of the Coast, 2003).
    \bibitem{ref46}J. K. Rowling, Harry Potter and the Half-Blood Prince (New York: Arthur A. Levine, 2005).
    \bibitem{ref47}Susan Caringella, "Affirmative Consent Reform Models," Addressing Rape Reform in Law and Practice (New York: Columbia University Press, 2009).
    \bibitem{ref48}Ibid.
    \bibitem{ref49}Walter Sinnot-Armstrong, "Consequentialism," Stanford Encyclopedia of Philosophy \url{http://plato.stanford.edu/entries/consequentialism/}
    \bibitem{ref50}Robert Kirk, "Zombies," The Stanford Encyclopedia of Philosophy \url{http://plato.stanford.edu/entries/zombies/}. See also Dennet's counter-arguments to the zombie concept in Daniel Dennett, Consciousness Explained (Boston: Little, Brown, 1991).
    \bibitem{ref51}Tanya Krzywinska, "Zombies in Gamespace: Form, Context, and Meaning in Zombie-Based Video Games," Zombie Culture: Autopsies of the Living Dead, Shawn McIntosh and Marc Leverette (ed.) (Lanham, MD: The Scarecrow Press, 2008).
    \bibitem{ref52}Postcolonial theorist Frantz Fanon posits that the zombie is an expression of native populations' feelings in response to the depersonalizing and confining pressures imposed by colonizing powers. Frantz Fanon, The Wretched of the Earth, Richard Philcox (trans.) (New York: Grove Press, 2004).
    \bibitem{ref53}Senko K. Maynard, "Japanese Communication in Global Context," Japanese Communication: Language and Thought in Context (Honolulu: University of Hawai'i Press, 1997).
    \bibitem{ref54}A style common, but hardly unique, to Japan, where it has a long history in the form of a rhetorical and narrative structure called kishotenketsu. This in turn is based on a type of four-line Chinese poem, in which the first line (ki) introduces a topic, the second (sho) develops it, the third (ten) introduces a surprise turn or apparent non sequitur, and the fourth (ketsu) integrates that surprise into the original topic in order to arrive at a novel or unexpected conclusion. See Maynard, "Japanese Communication Strategies: Collaboration toward Persuasion," Japanese Communication.
    \bibitem{ref55}Compare American Psychiatric Association, Diagnostic and Statistical Manual of Mental Disorders, Fourth Edition, Text Revision: DSM-IV-TR (Washington, DC: American Psychiatric Publishing, 2000).
    \bibitem{ref56}Hans Christian Anderson, "The Little Mermaid," The Classic Fairy Tales, Maria Tatar (ed.) (New York: W.W. Norton and Company, 1999). Excerpted from Anderson, Eighty Fairy Tales, R. P. Keigwin (trans.) (Pantheon Books, 1976).
    \bibitem{ref57}Ana Mardoll, "The Little Mermaid," Ramblings. \url{http://www.anamardoll.com/2012/05/disney-little-mermaid.html}
    \bibitem{ref58}American Psychiatric Association, DSM-IV-TR.
    \bibitem{ref59}Sharalyn Orbaugh, "Busty Battlin' Babes."
    \bibitem{ref60}Bewitched: The Complete Series (Sony Pictures Home Entertainment, 2013)
    \bibitem{ref61}Sally the Witch (Toei Animation, 1966)
    \bibitem{ref62}Sandra M. Gilbert and Susan Gubar, "[Snow White and Her Wicked Stepmother]," The Classic Fairy Tales. Excerpted from Gilbert and Gubar, The Madwoman in the Attic: The Woman Writer and the Nineteenth-Century Literary Imagination (New Haven: Yale University Press, 1979).
    \bibitem{ref63}Kumiko Fujimura-Fanselow, "The Japanese Ideology of 'Good Wives and Wise Mothers,': Trends in Contemporary Research," Gender and History, Volume 3, Issue 3 (1991).
    \bibitem{ref64}Smyers, "Symbolizing Inari: The Jewel."
    \bibitem{ref65}Noriko T. Reider, "Yamauba, the Mountain Ogress: Old Hag to Voluptuous Mother," Japanese Demon Lore: Oni from Ancient Times to the Present (Logan, Utah: Utah State University Press, 2010).
    \bibitem{ref66}Though at least this last may be merciful, if she chooses. Lafcadio Hearns, "Yuki-Onna," Kwaidan: Stories and Studies of Strange Things (Project Gutenberg, 1998). \url{http://www.gutenberg.org/files/1210/1210-h/1210-h.htm#yukionna}
    \bibitem{ref67}Gilbert and Gubar, "[Snow White and Her Wicked Stepmother]."
    \bibitem{ref68}Revolutionary Girl Utena: The Student Council Saga (Right Stuf, 2011).
    \bibitem{ref69}Revolutionary Girl Utena: The Apocalypse Saga (Right Stuf, 2011).
    \bibitem{ref70}For a detailed exploration of menstrual taboos and misogyny, see Janice Delaney, Mary Jane Lupton, and Emily Toth, The Curse: A Cultural History of Menstruation (Urbana, Ill.: University of Illinois Press, 1988 (Revised)).
    \bibitem{ref71}Saito, Beautiful Fighting Girl.
    \bibitem{ref72}Patrick W. Galbraith, "Moe: Exploring Virtual Potential in Post-Millennial Japan," Electronic Journal of Contemporary Japanese Studies, Volume 9, Issue 3 (October-November 2009). Note in particular the section titled "Otaku discussions of moe."
    \bibitem{ref73}Neera K. Badhwar and Roderick T. Long, "Ayn Rand," Stanford Encyclopedia of Philosophy \url{http://plato.stanford.edu/entries/ayn-rand/}
    \bibitem{ref74}Urobuchi, Fate/zero Volume 1.
    \bibitem{ref75}Smart, "Buddhism," The World's Religions (Cambridge, UK: Cambridge University Press, 1998 (2nd ed.)).
    \bibitem{ref76}Smart, "The Formation of Zen," The World's Religions.
    \bibitem{ref77}Vaclav Havel, "The Politics of Hope," Disturbing the Peace: A Conversation with Karel Hvizdala, Paul Wilson (trans.) (New York: Alfred A. Knopf, 1990).
    \bibitem{ref78}The wheel of dharma, fate, is a classic symbol of Buddhism. See Smart, "Buddhism."
    \bibitem{ref79}This is, generally speaking, how reincarnation works in both Hinduism and Buddhism—no physical or mnemonic trace of the previous life is maintained, but its karmic burden endures. See Smart, "The Ingredients of Indian Religion," The World's Religions.
    \bibitem{ref80}Smart, "Buddhism."
    \bibitem{ref81}Smart, "The Formation of Zen," The World's Religions.
    \bibitem{ref82}Adrian Stokes, Michelangelo: A Study in the Nature of Art (London: Tavistock Publications, 1955).
    \bibitem{ref83}Frank Lynn Meshberger, "An Interpretation of Michelangelo's Creation of Adam Based on Neuroanatomy," JAMA: Journal of the American Medical Association, Volume 264, Issue 14 (October 10, 1990).
    \bibitem{ref84}Note in particular item 1 in the ICD-10 description of dissocial personality disorder, "Callous unconcern for the feelings of others." World Health Organization, International Statistical Classification of Diseases and Related Health Problems, 10th edition (World Health Organization, 1994).
    \bibitem{ref85}On the eve of the festival of St. Walpurga (May 1), Walpurgisnacht, German folklore holds that witches and demons meet and dance on Mt. Brocken. See John Michael Cooper, "The Cultural and Religious Prehistories," Mendelssohn, Goethe, and the Walpurgis Night: The Heathen Muse in European Culture, 1700-1850 (Rochester, NY: University of Rochester Press, 2007).
    \bibitem{ref86}Giacomo Oreglia, The Commedia Dell'Arte, trans. Lovett F. Edwards (New York: Octagon Books, 1982).
    \bibitem{ref87}Smart, "Buddhism."
    \bibitem{ref88}A saying attributed to Lin-Chi Yihuan, an influential Tang Dynasty monk. See Burton Watson, The Zen Teachings of Master Lin-Chi: A translation of the Lin-chi lu (New York: Columbia University Press, 1999)
    \bibitem{ref89}Johann Wolfgang von Goethe, Faust, trans. Bayard Taylor (Hazleton, Penn.: Electronic Classic Series, 2005). \url{http://www2.hn.psu.edu/faculty/jmanis/goethe/goethe-faust.pdf}. This PDF, produced by Pennsylvania State University, covers only the first part of Faust.
    \bibitem{ref90}The Pure Land schools of Buddhism, broadly and briefly, teach of the Buddha Amida, who created a "pure land" to which the faithful will ascend upon death, and from which it is much easier to attain nirvana. Smart, "The Pure Land Movement and the Mappo," The World's Religions.
    \bibitem{ref91}Kanon being an important female Bodhisattva. See Smart, "The Pure Land Movements and the Mappo."
    \bibitem{ref92}"Sakura and Her True Feelings," Cardcaptor Sakura.
    \bibitem{ref93}Urobuchi, Fate/zero Volume 1.
    \bibitem{ref94}Goethe, Faust. \url{http://www2.hn.psu.edu/faculty/jmanis/goethe/goethe-faust.pdf}. Unless otherwise indicated, all references to the plot of Faust in this section refer to this version.
    \bibitem{ref95}Calvin Thomas, "Introduction to Faust," The German Classics of the Nineteenth and Twentieth Centuries Volume 1, Kuno Francke (ed.) (Project Gutenberg, 2004). Accessed via \url{http://www.gutenberg.org/ebooks/11123.}
    \bibitem{ref96}Elena P. O'Brien, "The Meaning of the Eternal Feminine in Goethe's Faust," Electronic Theses, Treatises, and Dissertations (2012). Paper 5076. \url{http://diginole.lib.fsu.edu/etd/5076/}
    \bibitem{ref97}Smart, "Buddhism."
    \bibitem{ref98}Ibid.
    \bibitem{ref99}Smart, "The Ingredients of Indian Religion."
    \bibitem{ref100}Most particularly in one of the sacred texts of Mahayana Buddhism, the Lotus Sutra. See Smart, "Nichiren and Japanese Nationalism," The World's Religions.
    \bibitem{ref101}Smart, "Buddhism."
    \bibitem{ref102}Smart notes the similarity between the figures of Kuanyin and Mary in Buddhism and Christianity, for example. See Smart, Buddhism and Christianity.
    \bibitem{ref103}Gilbert Gubar, "[Snow White and Her Wicked Stepmother]."
    \bibitem{ref104}The Turkey City Lexicon is a glossary of critical terms specific to science fiction, developed by and circulated to writers' workshops. Lewis Shiner and Bruce Sterling, Turkey City Lexicon—A Primer for SF Workshops (Science Fiction Writers of America, 2009). \url{http://www.sfwa.org/2009/06/turkey-city-lexicon-a-primer-for-sf-workshops/}
    \bibitem{ref105}"Puella Magi Oriko Magica," Puella Magi Wiki (2015). \url{http://wiki.puella-magi.net/Puella_Magi_Oriko_Magica}
    \bibitem{ref106}American Psychiatric Association, DSM-IV-TR.
    \bibitem{ref107}Ibid.
    \bibitem{ref108}Smart, "Buddhism."
    \bibitem{ref109}This is of course an extremely simplistic view that conflates entropy with the second law of thermodynamics, but will suffice for discussion as it is the view that Madoka Magica itself takes. And while simple, it is not inaccurate; see Martin Goldstein and Inga F. Goldstein, The Refrigerator and the Universe: Understanding the Laws of Energy (Cambridge, Mass: Harvard University Press, 1993).
    \bibitem{ref110}Smart, "The Ingredients of Indian Religion."
    \bibitem{ref111}Paracelsus, On the Nature of Things.
    \bibitem{ref112}Gerald Klingaman, "Plant of the Week: Red Spiderlily," University of Arkansas Division of Agriculture, Cooperative Extension Services (2000). Archived at \url{http://www.webcitation.org/61fHvTrxK.}
    \bibitem{ref113}Airplane! (Amazon Instant Video). \url{http://www.amazon.com/Airplane-Kareem-Abdul-Jabbar/dp/ B001K37CT8/ref=tmm_aiv_title_0?_encoding=UTF8&sr=&qid=}
    \bibitem{ref114}Iron Man (Amazon Instant Video). \url{http://www.amazon.com/Iron-Man-Robert-Downey-Jr/dp/B001KZG99A/ref=sr_1_2?ie=UTF8 &qid= 1419703982&sr=8-2}
    \bibitem{ref115}See for instance this image from the influential Rider-Waite tarot deck: \url{http://en.wikipedia.org/wiki/The_Fool_%28Tarot_card%29#mediaviewer/File:RWS_Tarot_00_Fool.jpg}
    \bibitem{ref116}Siegfried Gohr, Magritte: Attempting the Impossible (New York: Distributed Art Publishers, 2009)
    \bibitem{ref117}Wim Tigges, "The Themes and Motifs of Nonsense," The Anatomy of Literary Nonsense (Amsterdam: Rodopi, 1988)
    \bibitem{ref118}Alison Reike discusses this "category mistake" as one of four major types of nonsensical language, of which the other three are more straightforwardly wordplay, while the category mistake can appear on both linguistic and narrative levels. Reike, "Secret Sense, Coherent Nonsense," The Senses of Nonsense (Iowa City, Iowa: University of Iowa Press, 1992).
    \bibitem{ref119}Fujimura-Fanselow, "The Japanese Ideology of 'Good Wives and Wise Mothers.'"
    \bibitem{ref120}Arthur Edward Waite, The Book of Black Magic and Ceremonial Magic (revised) (San Diego, California: The Book Tree, 2006).
    \bibitem{ref121}Ibid.
    \bibitem{ref122}Isaiah 14:12-15, New International Version (Colorado Springs: Biblica, 2011).
    \bibitem{ref123}Smart, "The Nature of a Religion."
    \bibitem{ref124}Roy Flannagan, "Paradise Lost: Introduction," The Riverside Milton (Boston: Houghton Mifflin Company, 1998).
    \bibitem{ref125}William Blake, The Marriage of Heaven and Hell (reprint) (New York: Dover Books, 1994)
    \bibitem{ref126}John Milton, Paradise Lost. Reprinted in and annotated by Flannagan, The Riverside Milton.
    \bibitem{ref127}All references to the story of "The Nutcracker" in this chapter and the next refer to E.T.A Hoffmann, "The Nutcracker and the Mouse King," L.R.C. (trans.) \url{http://www.springhole.net/writing/the_nutcracker_and_the_mouse_king/index.html}
    \bibitem{ref128}Magica Quartet, Puella Magi Madoka Magica Production Note.    
\end{thebibliography}

\paragraph{About the Author}~{}

Jed A. Blue is a third-generation geek and lifelong animation buff. He has a degree in English from George Mason University, and lives in Washington, D.C., where he works as a technical writer. He can frequently be found giving panels on animation, literary and film criticism, and related topics at fan conventions along the East Coast. His favorite meguca is Homura, his favorite captain is the Sisko, and his favorite Doctor is Sylvester McCoy.

He blogs at JedABlue.com, where he is currently engaged in a multi-year critical study of the DC Animated Universe.
